\section{Forecast contract and controls}
Let $V_t$ be the entities active after processing all events recorded by origin
$t$, and let $g_r(i,t)$ denote entity $i$'s observed membership in relation type
$r$. For ordered horizons $h_1<\cdots<h_K$, the outcome
$Y_{ik}=\ind\{T_i\leq t+h_k\}$ records the first exit from the current task
state or the first qualifying application resolution. Each forecast draw
assigns one of $K+1$ categories, with the final category indicating survival
beyond $h_K$. This contract does not predict subsequent returns, exact
within-bin times, or outcomes of future arrivals.

An individual predictor supplies monotone CDFs $F_{ik}$. For a current group
$A$, the expected exit count is $\sum_{i\in A}F_{ik}$ under every coupling,
whereas its variance also includes pairwise outcome covariances. Thus a
coupling can change aggregate dispersion and tails but cannot fix a biased
expected count. We exploit this distinction in two separate experiments:
matched-marginal dependence comparisons and controlled marginal correction.

\subsection{Entity-coherent factor scenarios}
For each draw, generate independent standard normal global, relation-group,
and private factors and set
\begin{equation}
 Z_i=a_iZ_0+\sum_r b_{ir}Z_{r,g_r(i,t)}+c_i\epsilon_i,
 \quad a_i^2+\sum_r b_{ir}^2+c_i^2=1.
\end{equation}
Using $U_i=\Phi(Z_i)$ and $Y_{ik}=\ind\{U_i\leq F_{ik}\}$ preserves each
individual CDF and gives ordered horizon outcomes. Shared draws persist across
streamed entity chunks. Missing memberships transfer their squared loading
to private noise. Identities index incidence; no identity embedding or future
assignment is used. The construction handles simultaneous job and machine
membership without a dense entity covariance matrix.

Independent sampling uses only private noise. Global models learn
state-specific global/private weights; typed models add one weight per
observed relation type. The initial study also includes context-dependent
weight offsets and an exponential common-shock family. All weights are
nonnegative, with at least .02 private variance. These families cannot express
arbitrary negative association. We fit them by composite binary pair
likelihood using earlier out-of-sample individual forecasts, random fitting
horizon pairs, and January or original-development checkpoint selection.
Different calibration treatments receive separate dependence fits in the
fresh replication; retrospective BPI 2017 diagnostics retain original weights
to isolate marginal correction.

We permute noise-sharing incidence within current state and half-unit
last-horizon logit-risk bins, retaining true evaluation groups. This preserves
individual probabilities, fitted weights and group sizes within strata.
Additional Borg controls remove job or machine variance and transfer it to
private noise. These sampler interventions assess predictive grouping under a
fitted model, not causal effects of reassigning actual tasks or people.

\subsection{Delayed calibration}
A shared logit intercept preserves horizon order:
\begin{equation}
 F_{ik}^{(\beta)}=\operatorname{logit}^{-1}
 \{\operatorname{logit}(F_{ik})+\beta\}.
\end{equation}
We fit $\beta\in[-2,2]$ by equal-origin, equal-horizon Bernoulli loss with
$10^{-4}\beta^2$ regularization. The static control uses an earlier calibration
window. Rolling correction uses origins in the previous 60 days whose entire
14-day outcome window has matured; fewer than four eligible origins triggers
an expanding-history fallback. Updated labels enter only after maturity.

A retrospective failure motivates an explicit version control: calibration
residuals from earlier separately fitted predictors need not describe the
current predictor's bias. We therefore compare that transfer with history
from the same fitted predictor. This is a practical calibration control,
not a new online-learning algorithm. Static calibration data also selected
the original neural checkpoint/temperature, so its own fit score is not an
independent validation estimate.

\subsection{Direct counts and discrete reconciliation}
For group $A$, average the individual first-exit category probabilities to
obtain $\bar{\pi}_A$. A conditional Dirichlet--multinomial (DM) baseline uses
\begin{equation}
 N_A\sim\operatorname{DM}(|A|,\kappa q_A),\qquad
 q_A=\operatorname{softmax}(\log\bar\pi_A+\delta).
\end{equation}
We fit one concentration and optionally four category offsets (the fifth is
zero) per projection type, with equal-origin and equal-within-origin-group
likelihood weights. Offsets have fixed weak shrinkage. Concentration-only
and bias-adjusted candidates are selected on development count CRPS.
Training uses complete out-of-sample cohorts, never sampled partial group
counts. Every group draw is bounded, integer-valued and ordered across
horizons. Separately drawn overlapping projections need not correspond to a
common entity realization.

For the retrospective original-log analyses we also fit DM on the more recent
development forecasts, matching the fitting window of historical residuals.
Leave-one-origin-out development CRPS selects concentration-only versus
bias-adjusted DM; all development origins then fit the final predictor.
These selection folds can use later development origins and are not online
validation. Exact beta-binomial cumulative-count marginals supply CRPS and
discrete quantiles for this control; joint-profile scores still use sampled
DM trajectories. This stronger comparison avoids attributing an advantage
to dependence when it comes from access to recent same-version errors.

For the workflow hierarchy, include every handler group, plus an unknown
handler group, in a complete partition. Independently sampled bottom count
vectors $N_g$ are reweighted by the population forecast mass
\begin{equation}
 w^{(m)}\ \propto\ p_{\rm pop}\!\left(\sum_gN_g^{(m)}\right).
\end{equation}
This two-level conditioning adaptation \citep{zambon2024} uses 8,192 proposal
draws and resamples 512 coherent group/population trajectories. We report
effective sample sizes and verify its distribution against an enumerated
small example. Bottom-up aggregation without conditioning is a separate
control. Neither count method preserves individual CDFs. The implemented
hierarchy is complete for workflow handlers; it is not a full crossed
job--machine reconciliation implementation. Historical aggregate residual
resampling remains a comparator in the original two logs.

\section{Evaluation}
We score current-state population counts and observed relation groups of at
least five entities. Fair ensemble CRPS \citep{ferro2014} is divided by current
group size, then averaged equally over groups within projection, horizons,
projection types and origins. This weighting is a benchmark convention, not
an inferred organizational utility. Coverage and width of 90\% intervals are
reported separately by projection.

For remaining current-cohort count $R=|A|-C$, the reserve proxy chooses forecast
quantile $b$ at $r/(r+1)$ and scores
\begin{equation}
 \ell_r(b,R)=\{r(R-b)_++(b-R)_+\}/|A|,\quad r\in\{3,9\}.
\end{equation}
It measures a specified quantile tradeoff, not money or realized staffing
savings. CRPS integrates quantile losses \citep{gneiting2007}; improving it
does not guarantee improvement at this particular quantile.

For joint group risk, let $x_g=C_g/|g|$. We use an energy score with
$d(x,y)=\sqrt{G^{-1}\sum_g(x_g-y_g)^2}$ and a power-.5 variogram score
\citep{scheuerer2015}. The latter averages 2,048 outcome-independent ordered
group pairs; subtracting the ensemble variance divided by draw count removes
finite-ensemble bias of its squared moment error. Its estimate can be slightly
negative. We additionally score the distribution of
\begin{equation}
 D_k=\max_g\frac{\mu_{gk}-C_{gk}}
 {\sqrt{1+\sum_{i\in g}F_{ik}(1-F_{ik})}},
 \qquad \mu_{gk}=\sum_{i\in g}F_{ik}.
\end{equation}
This is the largest standardized slow-progress deviation relative to a
fixed reference forecast. Every compared method uses the same reference
$F$, including recalibrated methods. It measures forecast-relative group
risk, not an observed capacity threshold or an organizational SLA.

We average three random seeds within each origin before paired circular
moving-block resampling (10,000 replicates). Borg blocks span one/two days;
BPI 2017 blocks span 14/28 days; the fresh BPI 2012 analysis prespecifies
seven days with fourteen-day sensitivity. Short traces provide few blocks,
so intervals describe these observed periods and offer limited evidence for
other operating regimes. Cases and random seeds are not independent temporal
replications.
